\section{A rank measure from differentiation and multiplication}
\label{sec:rank}

We associate a linear map to a polynomial by differentiating in one
group of variables and multiplying in a disjoint group.  The rank of
this map is subadditive.  More importantly, the map associated to a
product can be factored through intermediate homogeneous spaces, whose
dimensions will control the circuit upper bound.

\subsection{The operator and its homogeneous spaces}

Let $K$ be a field, and let $V,U$ be disjoint finite sets of variables,
of cardinalities $v,u\ge1$.  For a variable set $E$ and an integer
$t\ge0$, write $\HH_E(t)=K[E]_t$ for the vector space of homogeneous
polynomials of degree $t$.  Its dimension is
\[
 h(|E|,t),\qquad h(r,t)=\binom{r+t-1}{t}.
\]
For $Q\in K[V,U]$, let
\[
 T_Q=Q(\partial_V,U)
\]
be the operator on $K[V,U]$ obtained by replacing each variable in $V$
by its formal partial derivative and each variable in $U$ by
multiplication by that variable.  These operators commute because the
two variable sets are disjoint.  Thus $T_{Q_1Q_2}=T_{Q_1}T_{Q_2}$.
This identity holds over every field.

Fix integers $k,m\ge0$, $a\ge k$, and $b\ge0$.  If $g$ is homogeneous
of degree $k+m$, denote its component of $V$-degree $k$ and $U$-degree
$m$ by $g_{[k,m]}$.  Define the finite-dimensional map
\begin{equation}
\label{eq:rank-map}
 F_g=g_{[k,m]}(\partial_V,U):
 \HH_V(a)\otimes\HH_U(b)
 \longrightarrow
 \HH_V(a-k)\otimes\HH_U(b+m),
 \qquad R(g)=\rank F_g.
\end{equation}
The source dimension is $D=h(v,a)h(u,b)$.  Projection to a bidegree
and operator substitution are both linear, so
\[
 R(g_1+g_2)\le R(g_1)+R(g_2),\qquad
 R(cg)=R(g)\quad(c\in K\setminus\{0\}).
\]
Also $R(0)=0$.  The multiplicative identity for $T_Q$ does not assert
multiplicativity of $g\mapsto F_g$: the latter includes a fixed
bidegree projection.  We will first decompose products into their
bidegree components and then use multiplicativity of $T_Q$ on each
contribution.

\subsection{Choosing the degrees}

We now choose the spaces in~\eqref{eq:rank-map} for degree-$n$
polynomials on the $n$ matrix layers.  Throughout the proof $n$ is
sufficiently large, and all implicit constants and thresholds are
absolute.  Let $s$ be the largest integer at most $\sqrt n$ having the
same parity as $n$, and put
\begin{equation}
\label{eq:parameters}
\begin{gathered}
 k=\frac{n-s}{2},\qquad m=\frac{n+s}{2},\qquad
 \rho=\frac{k}{m},\qquad \lambda=\frac{k}{n}
       =\frac{\rho}{1+\rho},\\
 v=kn^2,\qquad u=mn^2,\qquad
 a=\left\lceil\frac{v}{\sqrt n-1}\right\rceil,
 \qquad \alpha=\frac{a}{a+v},\\
 b=\left\lceil\frac{u\alpha^\rho}{1-\alpha^\rho}\right\rceil,
 \qquad \beta=\frac{b}{b+u},\qquad
 q=\alpha^{(1+\rho)/2},\qquad D=h(v,a)h(u,b).
\end{gathered}
\end{equation}
For now assign any $k$ full matrix layers to $V$ and the other $m$
layers to $U$.  Each layer retains all its $n^2$ variables, including
variables absent from the endpoint entry of the matrix product.
Section~\ref{sec:moments} will specify their order when bounding the
rank of $\IMM_{n,n}$ from below.

The small excess $m-k=s$ makes the slope $\lambda$ slightly less than
$1/2$.  For small positive even factor degrees $e$, this moves
$\lambda e$ away from the integers; Section~\ref{sec:circuits}
quantifies the resulting rank saving.  The definitions of $a,b$
balance the dimension decrease caused by differentiation against the
dimension increase caused by multiplication.  Without the ceiling in
the definition of $b$, we would have $\beta=\alpha^\rho$ and
$\alpha^k\beta^{-m}=1$.  We record the rounding estimates explicitly,
as their errors must stay small after many factors are combined.

\begin{lemma}
\label{lem:parameters}
For all sufficiently large $n$, the parameters in~\eqref{eq:parameters}
satisfy
\begin{gather*}
 \frac{\sqrt n}{2}\le s\le\sqrt n,\qquad
 \lambda\ge\frac38,\qquad 0<\rho<1,\qquad
 v,u=\Theta(n^3),\qquad a,b=\Theta(n^{5/2}),\\
 a\ge k,\qquad 2n\le\frac12\min(a,b),\qquad
 n^{-1/2}\le\alpha=n^{-1/2}\exp(O(1/a)),\qquad
 \alpha^{-1}\le\sqrt n,\\
 \alpha^\rho\le\beta\le\alpha^\rho\exp(C/b),
 \qquad \beta\le\frac2{\sqrt n},\\
 q=n^{-1/2}\exp\bigl(O(s\log n/n+1/a)\bigr),
 \qquad n^{-1/2}\le q\le\frac2{\sqrt n}
 \le n^{-2/5},\qquad q\le\frac12,
\end{gather*}
where $C$ is an absolute constant.
\end{lemma}

\begin{proof}
The parity choice gives $\sqrt n-2<s\le\sqrt n$, and hence the
claims about $s,k,m,\lambda,\rho,v,u$.  Write
$a_0=v/(\sqrt n-1)$, so that $a_0\le a<a_0+1$ and
$a_0/(a_0+v)=n^{-1/2}$.  The function $x/(x+v)$ is increasing, and
the derivative of its logarithm is $v/(x(x+v))\le1/x$.
The mean value theorem therefore gives
\[
 0\le\log(\alpha n^{1/2})\le\frac1{a_0}=O(1/a).
\]
In particular $a=\Theta(n^{5/2})$ and $\alpha^{-1}\le\sqrt n$.

Since $1-\rho=s/m=O(s/n)$, the preceding estimate implies
\[
 \alpha^\rho
 =n^{-1/2}\exp\bigl(O(s\log n/n+1/a)\bigr)
 =n^{-1/2}(1+o(1)).
\]
Put $b_0=u\alpha^\rho/(1-\alpha^\rho)$.  Then
$b_0=\Theta(n^{5/2})$, $b_0\le b<b_0+1$, and
$b_0/(b_0+u)=\alpha^\rho$.  Applying the same logarithmic derivative
bound with $u$ in place of $v$ gives
\[
 0\le\log(\beta/\alpha^\rho)\le1/b_0\le C/b.
\]
This proves the assertions about $b,\beta$, including
$\beta\le2/\sqrt n$ eventually.  The growth of $a,b$ also gives
$a\ge k$ and $2n\le\min(a,b)/2$.

Finally,
\[
 \log q=-\frac{1+\rho}{4}\log n+O(1/a)
 =-\frac12\log n+O(s\log n/n+1/a).
\]
Because $0<\alpha<1$ and $(1+\rho)/2\le1$, we have
$q\ge\alpha\ge n^{-1/2}$.  The remaining upper bounds follow from
the last display and $s\log n/n\to0$.
\end{proof}

The estimate $q\le n^{-2/5}$ will suffice for the main power savings.
The sharper $q=O(n^{-1/2})$ will be essential when summing the
smaller errors contributed by even factor degrees.
